\documentclass[11pt,a4paper]{article}
\usepackage[utf8]{inputenc}
\usepackage[vietnamese]{babel}
\usepackage{amsmath,amssymb}
\usepackage{geometry}
\geometry{a4paper, total={180mm,260mm}, left=15mm, right=15mm, top=18mm, bottom=18mm}

\usepackage{tikz}
\usetikzlibrary{shapes.geometric, arrows.meta, positioning, calc}
\tikzset{
  arr/.style={-{Stealth[scale=0.8]}},
  every node/.append style={align=center}
}
\usepackage{tcolorbox}
\tcbuselibrary{skins,breakable}
\usepackage{fontawesome}
\usepackage{fancyhdr}
\usepackage{enumitem}

\pagestyle{fancy}
\fancyhf{}
\renewcommand{\headrulewidth}{0.5pt}
\renewcommand{\footrulewidth}{0.5pt}
\fancyhead[L]{\small\bfseries\color{violet!85!black}\faBook\ SỔ TAY KIẾN THỨC TOÁN 2}
\fancyhead[R]{\small\bfseries\color{violet!85!black}Lớp Toán Thầy Lĩnh -- loptoanthaylinh.vn}
\fancyfoot[C]{\small\color{gray}-- Sổ tay kiến thức: Bài 47. Luyện tập chung hình khối --}

\setlength{\parindent}{0pt}
\setlength{\parskip}{0pt}

\begin{document}

% --- BANNER BÀI HỌC ---
\begin{tcolorbox}[
  enhanced,
  colback=violet!8!white,
  colframe=violet!70!black,
  arc=3mm, boxrule=1.2pt,
  left=3mm, right=3mm, top=2mm, bottom=2mm,
  halign=center
]
  {\LARGE\bfseries\color{violet!90!black}\faBookmark\ BÀI 47: LUYỆN TẬP CHUNG HÌNH KHỐI}\\[1mm]
  {\itshape\footnotesize\color{black!75} Đại tổng kết 4 khối không gian lớp 2: Khối lập phương, Khối hộp chữ nhật, Khối trụ và Khối cầu trên 1 trang A4}
\end{tcolorbox}

\vspace{-2mm}

\noindent
\begin{minipage}[t]{0.49\textwidth}\vspace{0pt}
\begin{tcolorbox}[
  enhanced,
  colback=violet!4!white,
  colframe=violet!75!black,
  arc=2mm, boxrule=0.8pt,
  left=2mm, right=2mm, top=1.8mm, bottom=1.8mm,
  title={\footnotesize\bfseries\color{white}\faKey\ 1. Bộ tứ hình khối lớp 2}
]
{\footnotesize
\begin{itemize}\setlength{\itemsep}{1pt}\setlength{\leftskip}{-2.5mm}
  \item \textbf{1. Khối lập phương:}\\
  Có $6$ mặt đều là các hình vuông bằng nhau. (VD: con súc sắc, khối rubik).
  \item \textbf{2. Khối hộp chữ nhật:}\\
  Có $6$ mặt là các hình chữ nhật. (VD: bao diêm, hộp bánh, viên gạch).
  \item \textbf{3. Khối trụ:}\\
  Có 2 đáy hình tròn phẳng, lăn được khi nằm ngang. (VD: lon sữa, cọc tiêu).
  \item \textbf{4. Khối cầu:}\\
  Mặt cong toàn phần, lăn mọi hướng. (VD: quả địa cầu, bóng đá).
\end{itemize}
}
\end{tcolorbox}
\end{minipage}
\hfill
\begin{minipage}[t]{0.49\textwidth}\vspace{0pt}
\begin{tcolorbox}[
  enhanced,
  colback=teal!4!white,
  colframe=teal!75!black,
  arc=2mm, boxrule=0.8pt,
  left=2mm, right=2mm, top=1.8mm, bottom=1.8mm,
  title={\footnotesize\bfseries\color{white}\faSliders\ 2. Kĩ năng đếm hình khối xếp chồng}
]
{\footnotesize
\begin{itemize}\setlength{\itemsep}{1pt}\setlength{\leftskip}{-2.5mm}
  \item Khi các khối lập phương nhỏ xếp chồng lên nhau thành mô hình:
  \item \textbf{Quy tắc đếm theo tầng (lớp):}\\
  - Đếm từ tầng dưới cùng (tầng 1) lên tầng trên.\\
  - Chú ý những khối \textbf{bị che khuất} ở phía dưới các khối tầng trên (khối phía trên muốn đứng được thì bên dưới bắt buộc phải có khối đỡ).
  \item \textbf{Xếp hình sáng tạo:} Ghép các khối để tạo thành lâu đài, ngôi nhà, cây cầu.
\end{itemize}
}
\end{tcolorbox}
\end{minipage}

\vspace{-1.5mm}
\begin{tcolorbox}[
  enhanced,
  colback=orange!4!white,
  colframe=orange!80!black,
  arc=2mm, boxrule=0.8pt,
  left=2mm, right=2mm, top=1.5mm, bottom=1.5mm,
  title={\footnotesize\bfseries\color{white}\faSitemap\ 3. Sơ đồ tư duy 4 nhánh: Phân loại 4 khối không gian cơ bản}
]
\centering
\begin{tikzpicture}[
  node distance=0.8cm,
  box/.style={rectangle, rounded corners=2pt, draw=violet!80!black, fill=violet!15, text width=2.4cm, align=center, font=\scriptsize\bfseries, inner sep=3pt},
  branch/.style={rectangle, rounded corners=2pt, draw=#1!80!black, fill=#1!10, text width=2.4cm, align=center, font=\scriptsize, inner sep=2.5pt},
  arr/.style={-{Stealth[scale=0.8]}, thick, #1!80!black}
]
  \node[box] (root) {HÌNH KHỐI\\KHÔNG GIAN};
  \node[branch=teal, above left=0.3cm and 0.4cm of root] (b1) {\textbf{KHỐI LẬP PHƯƠNG}\\$6$ mặt hình vuông\\Không lăn được};
  \node[branch=orange, above right=0.3cm and 0.4cm of root] (b2) {\textbf{HỘP CHỮ NHẬT}\\$6$ mặt hình chữ nhật\\Không lăn được};
  \node[branch=green, below left=0.3cm and 0.4cm of root] (b3) {\textbf{KHỐI TRỤ}\\$2$ đáy phẳng tròn\\Lăn ngang, trượt đứng};
  \node[branch=purple, below right=0.3cm and 0.4cm of root] (b4) {\textbf{KHỐI CẦU}\\Mặt cong toàn bộ\\Lăn về mọi hướng};

  \draw[arr=teal] (root) -- (b1);
  \draw[arr=orange] (root) -- (b2);
  \draw[arr=green] (root) -- (b3);
  \draw[arr=purple] (root) -- (b4);
\end{tikzpicture}
\end{tcolorbox}

\vspace{-1.5mm}
\noindent
\begin{minipage}[t]{0.49\textwidth}\vspace{0pt}
\begin{tcolorbox}[
  enhanced,
  colback=green!4!white,
  colframe=green!60!black,
  arc=2mm, boxrule=0.8pt,
  left=2mm, right=2mm, top=1.8mm, bottom=1.8mm,
  title={\footnotesize\bfseries\color{white}\faCheckCircle\ 4. Ví dụ thực hành chuẩn mực}
]
{\footnotesize
\textbf{Đếm khối lập phương nhỏ:}\\
Một hình gồm: Tầng dưới có $4$ khối, tầng trên có $2$ khối đặt lên trên.\\
$\implies$ Tổng số khối lập phương là: $4 + 2 = \mathbf{6}$ khối.
\textbf{Vật nào có dạng khối cầu:}\\
A. Hộp sữa bột \quad B. Quả bóng chuyền \quad C. Viên gạch.\\
$\implies$ Đáp án đúng là \textbf{B. Quả bóng chuyền}.
}
\end{tcolorbox}
\end{minipage}
\hfill
\begin{minipage}[t]{0.49\textwidth}\vspace{0pt}
\begin{tcolorbox}[
  enhanced,
  colback=red!4!white,
  colframe=red!75!black,
  arc=2mm, boxrule=0.8pt,
  left=2mm, right=2mm, top=1.8mm, bottom=1.8mm,
  title={\footnotesize\bfseries\color{white}\faExclamationTriangle\ 5. Mẹo tránh bẫy trắc nghiệm}
]
{\footnotesize
\begin{itemize}\setlength{\itemsep}{1pt}\setlength{\leftskip}{-2.5mm}
  \item \textbf{Bẫy đếm sót khối bị che khuất:}\\
  Nếu nhìn thấy một khối ở tầng 2, chắc chắn bên dưới nó phải có một khối ở tầng 1 làm trụ đỡ, không được bỏ quên khối này!
  \item Khối lập phương có tất cả các mặt đều bằng nhau.
\end{itemize}
}
\end{tcolorbox}
\end{minipage}


\end{document}
